\documentclass{article}
\usepackage{amsmath,amssymb,amsthm}
\usepackage{bm}
\usepackage{algorithm}
\usepackage{algorithmic}
\usepackage{hyperref}
\usepackage{geometry}
\geometry{margin=1in}
\newtheorem{theorem}{Theorem}
\newtheorem{lemma}{Lemma}
\newtheorem{assumption}{Assumption}
\newcommand{\E}{\mathbb{E}}
\newcommand{\R}{\mathbb{R}}

\begin{document}

\title{Differentially Private Stochastic Gradient Descent \\ (Gaussian‑Noised SGD)}
\author{}
\date{}
\maketitle

\tableofcontents
\newpage

%%%%%%%%%%%%%%%%%%%%%%%%%%%%%%%%%%%%%%%%%%%%%%%%%%%%%%%%%%%%%%%%%%%%%%%%%%%%%%%%
\section*{Algorithm Name}
\addcontentsline{toc}{section}{Algorithm Name}
\begin{center}
\Large \textbf{DP‑SGD with Gaussian Gradient Perturbation}
\end{center}
%%%%%%%%%%%%%%%%%%%%%%%%%%%%%%%%%%%%%%%%%%%%%%%%%%%%%%%%%%%%%%%%%%%%%%%%%%%%%%%%

\section*{Mathematical Formulation}
\addcontentsline{toc}{section}{Mathematical Formulation}
Let $f:\R^{d}\rightarrow\R$ be the (convex, $L$‑smooth) empirical risk
\[
f(w)=\frac{1}{n}\sum_{i=1}^{n}\ell(w;z_i),\qquad z_i\in\mathcal{Z}.
\]
At iteration $t$ we compute a (possibly minibatch) stochastic gradient
\[
\tilde g_t = \frac{1}{|B_t|}\sum_{i\in B_t}\nabla \ell(w_t;z_i)
\]
and then add isotropic Gaussian noise:
\[
g_t \;=\; \tilde g_t + \xi_t ,\qquad 
\xi_t\sim\mathcal N\!\bigl(0,\sigma^2 I_d\bigr).
\]
The parameter vector is updated by
\[
\boxed{w_{t+1}= w_t-\eta_t\, g_t }                                   \tag{1}
\]
where $\eta_t>0$ is the learning rate.  
The privacy accountant (e.g. moments accountant) tracks the cumulative
$(\varepsilon,\delta)$‑budget from the sequence of Gaussian mechanisms.

%%%%%%%%%%%%%%%%%%%%%%%%%%%%%%%%%%%%%%%%%%%%%%%%%%%%%%%%%%%%%%%%%%%%%%%%%%%%%%%%
\section*{Pseudocode}
\addcontentsline{toc}{section}{Pseudocode}
\begin{algorithm}[H]
\caption{DP‑SGD with Gaussian Noise}
\label{alg:dp-sgd}
\begin{algorithmic}[1]
\REQUIRE Initial point $w_0\in\R^{d}$, learning‑rate schedule $\{\eta_t\}_{t=0}^{T-1}$,
 minibatch size $b$, noise scale $\sigma$, total iterations $T$.
\FOR{$t=0$ \TO $T-1$}
    \STATE Sample a minibatch $B_t\subset\{1,\dots,n\}$ of size $b$ uniformly at random.
    \STATE Compute the (clipped) stochastic gradient 
    \[
    \tilde g_t=\frac{1}{b}\sum_{i\in B_t}\nabla\ell(w_t;z_i).
    \]
    \STATE Sample $\xi_t\sim\mathcal N(0,\sigma^2 I_d)$.
    \STATE Form the noisy gradient $g_t=\tilde g_t+\xi_t$.
    \STATE Update the parameters $w_{t+1}=w_t-\eta_t g_t$.
    \STATE (Optional) Update the privacy accountant with the Gaussian
            mechanism $(\sigma,b,n)$.
\ENDFOR
\RETURN $w_T$ (or an average $\bar w_T = \frac1T\sum_{t=0}^{T-1}w_t$).
\end{algorithmic}
\end{algorithm}
%%%%%%%%%%%%%%%%%%%%%%%%%%%%%%%%%%%%%%%%%%%%%%%%%%%%%%%%%%%%%%%%%%%%%%%%%%%%%%%%

\section*{Dry Run Example}
\addcontentsline{toc}{section}{Dry Run Example}
Consider the scalar quadratic
\[
f(w)=(w-3)^2,\qquad \nabla f(w)=2(w-3).
\]
Set $w_0=0$, constant step size $\eta=0.1$, noise variance $\sigma^2=0.01$,
and run three iterations.  For reproducibility we fix the random draws
$\xi_0=0.05$, $\xi_1=-0.02$, $\xi_2=0.03$.

\begin{center}
\begin{tabular}{c|c|c|c}
$t$ & $g_t=\nabla f(w_t)+\xi_t$ & $w_{t+1}=w_t-\eta g_t$ & $f(w_{t+1})$ \\ \hline
0 & $2(0-3)+0.05=-5.95$ & $0-0.1(-5.95)=0.595$ & $(0.595-3)^2=5.777$ \\
1 & $2(0.595-3)-0.02=-4.83$ & $0.595-0.1(-4.83)=1.078$ & $(1.078-3)^2=3.691$ \\
2 & $2(1.078-3)+0.03=-3.862$ & $1.078-0.1(-3.862)=1.464$ & $(1.464-3)^2=2.365$
\end{tabular}
\end{center}
The objective value decreases at each step despite the injected noise.

%%%%%%%%%%%%%%%%%%%%%%%%%%%%%%%%%%%%%%%%%%%%%%%%%%%%%%%%%%%%%%%%%%%%%%%%%%%%%%%%
\section*{Assumptions}
\addcontentsline{toc}{section}{Assumptions}
\begin{assumption}[Convexity]\label{ass:convex}
$f$ is convex: $\forall w,w',\; f(w')\ge f(w)+\langle\nabla f(w),w'-w\rangle$.
\end{assumption}
\begin{assumption}[Smoothness]\label{ass:smooth}
$f$ is $L$‑smooth: $\forall w,w',\;
\|\nabla f(w)-\nabla f(w')\|\le L\|w-w'\|$,
equivalently
\[
f(w')\le f(w)+\langle\nabla f(w),w'-w\rangle+\frac{L}{2}\|w'-w\|^2.
\]
\end{assumption}
\begin{assumption}[Bounded Domain]\label{ass:bounded}
There exists $D>0$ such that $\|w_t-w^\star\|\le D$ for all $t$, where
$w^\star\in\arg\min f$.
\end{assumption}
\begin{assumption}[Noise]\label{ass:noise}
The added noise is i.i.d. Gaussian,
$\xi_t\sim\mathcal N(0,\sigma^2 I_d)$, independent of $w_t$ and the minibatch.
Consequently $\E[\xi_t]=0$ and $\E[\|\xi_t\|^2]=d\sigma^2$.
\end{assumption}
\begin{assumption}[Unbiased Stochastic Gradient]\label{ass:unbiased}
$\E[\tilde g_t\mid w_t]=\nabla f(w_t)$ and $\E[\|\tilde g_t-\nabla f(w_t)\|^2]\le \zeta^2$
for some $\zeta\ge0$ (the usual SGD variance bound).
\end{assumption}

%%%%%%%%%%%%%%%%%%%%%%%%%%%%%%%%%%%%%%%%%%%%%%%%%%%%%%%%%%%%%%%%%%%%%%%%%%%%%%%%
\section*{Main Convergence Theorem}
\addcontentsline{toc}{section}{Main Convergence Theorem}
\begin{theorem}[Expected Suboptimality of DP‑SGD]\label{thm:main}
Assume \ref{ass:convex}–\ref{ass:noise} and choose a constant stepsize
$\eta\in(0,\tfrac{1}{L}]$. Let $w_T$ be the iterate after $T$ updates
according to \eqref{1}. Then
\[
\E\!\bigl[f(w_T)-f(w^\star)\bigr]
\;\le\;
\frac{\|w_0-w^\star\|^2}{2\eta T}
\;+\;
\frac{\eta L}{2}\bigl(\zeta^2+d\sigma^2\bigr).
\tag{2}
\]
In particular, setting $\eta=\Theta\!\bigl(1/\sqrt{T}\bigr)$ yields
\[
\E\!\bigl[f(w_T)-f(w^\star)\bigr]=O\!\Bigl(\frac{1}{\sqrt{T}}\Bigr)
\;+\;O\!\bigl(\sigma^2\sqrt{T}\bigr).
\]
\end{theorem}
\begin{theorem}[Averaged Iterate]\label{thm:avg}
If we output the uniform average $\bar w_T=\frac{1}{T}\sum_{t=0}^{T-1}w_t$, then
\[
\E\!\bigl[f(\bar w_T)-f(w^\star)\bigr]
\;\le\;
\frac{\|w_0-w^\star\|^2}{2\eta T}
\;+\;
\frac{\eta L}{2}\bigl(\zeta^2+d\sigma^2\bigr).
\tag{3}
\]
\end{theorem}
Both bounds match the classical SGD rate up to an additive term
proportional to the noise variance $d\sigma^2$.

%%%%%%%%%%%%%%%%%%%%%%%%%%%%%%%%%%%%%%%%%%%%%%%%%%%%%%%%%%%%%%%%%%%%%%%%%%%%%%%%
\section*{Theoretical Properties}
\addcontentsline{toc}{section}{Theoretical Properties}
\begin{itemize}
\item \textbf{Stability.} The bound \eqref{2} is linear in the noise variance.
When $\sigma\to0$ we recover the standard SGD rate.
\item \textbf{Hyper‑parameter sensitivity.}
The stepsize $\eta$ appears both in the $1/(\eta T)$ term and in the
$\eta L$ term. The optimal choice $\eta^\star = \sqrt{\frac{\|w_0-w^\star\|^2}
{L(\zeta^2+d\sigma^2)T}}$ balances the two contributions.
\item \textbf{Comparison to baselines.}
For pure SGD ($\sigma=0$) Theorem \ref{thm:main} reduces to the classic
$O(1/\sqrt{T})$ bound. Adding Gaussian noise degrades the rate by an
additive $O(\eta\sigma^2)$ term, which is unavoidable under differential
privacy (the noise magnitude is dictated by the privacy budget).
\item \textbf{Privacy‑utility trade‑off.}
The noise scale $\sigma$ is calibrated to $(\varepsilon,\delta)$ via the
Gaussian mechanism. Larger $\varepsilon$ (weaker privacy) permits smaller
$\sigma$, thus improving the bound.
\end{itemize}

\section*{Discussion}
\addcontentsline{toc}{section}{Discussion}
The theorem shows that DP‑SGD retains the $O(1/\sqrt{T})$ convergence
behaviour of vanilla SGD as long as the noise variance grows slower than
$1/T$. In practice one selects $\sigma$ from a privacy accountant,
chooses $\eta$ according to the optimal balance above, and often
averages iterates to reduce variance. The analysis extends verbatim to
mini‑batching (the variance term $\zeta^2$ then scales with $1/b$).

%%%%%%%%%%%%%%%%%%%%%%%%%%%%%%%%%%%%%%%%%%%%%%%%%%%%%%%%%%%%%%%%%%%%%%%%%%%%%%%%
\section*{Preliminaries (Definitions and Lemmas)}
\addcontentsline{toc}{section}{Preliminaries (Definitions and Lemmas)}
\begin{lemma}[Smoothness inequality]\label{lem:smooth}
If $f$ is $L$‑smooth then for any $w,u\in\R^d$,
\[
f(u)\le f(w)+\langle\nabla f(w),u-w\rangle+\frac{L}{2}\|u-w\|^2 .
\]
\end{lemma}
\begin{lemma}[Quadratic expansion]\label{lem:exp}
For any vectors $a,b\in\R^d$ and scalar $\eta>0$,
\[
\|a-\eta b\|^2 = \|a\|^2 -2\eta\langle a,b\rangle +\eta^2\|b\|^2 .
\]
\end{lemma}
Both lemmas are elementary and will be invoked repeatedly.

%%%%%%%%%%%%%%%%%%%%%%%%%%%%%%%%%%%%%%%%%%%%%%%%%%%%%%%%%%%%%%%%%%%%%%%%%%%%%%%%
\section*{Main Proof (Step‑by‑Step Derivation)}
\addcontentsline{toc}{section}{Main Proof (Step-by-Step Derivation)}
We prove Theorem \ref{thm:main}.  All expectations are taken w.r.t. the
randomness of the minibatch, the Gaussian noise and any internal
randomness of the algorithm.

\subsection*{Step 1 – Smoothness bound on the function value}
By Lemma~\ref{lem:smooth} with $w=w_t$ and $u=w_{t+1}=w_t-\eta g_t$,
\begin{align}
f(w_{t+1})
&\le f(w_t)+\langle\nabla f(w_t),w_{t+1}-w_t\rangle
      +\frac{L}{2}\|w_{t+1}-w_t\|^2                         \nonumber\\
&= f(w_t)-\eta\langle\nabla f(w_t),g_t\rangle
   +\frac{L\eta^2}{2}\|g_t\|^2 .                              \tag{4}
\end{align}
[Step 1]

\subsection*{Step 2 – Decompose the noisy gradient}
Recall $g_t=\tilde g_t+\xi_t$ and $\E[\xi_t]=0$.  Using linearity of inner
products,
\[
\langle\nabla f(w_t),g_t\rangle
= \langle\nabla f(w_t),\tilde g_t\rangle
  +\langle\nabla f(w_t),\xi_t\rangle .                       \tag{5}
\]
Taking expectation conditioned on $w_t$ and the minibatch,
the second term vanishes because $\xi_t$ is independent and zero‑mean:
\[
\E\!\bigl[\langle\nabla f(w_t),\xi_t\rangle\mid w_t\bigr]=0 .
\] 
Moreover, by Assumption~\ref{ass:unbiased},
\[
\E\!\bigl[\tilde g_t\mid w_t\bigr]=\nabla f(w_t).
\] 
Hence
\[
\E\!\bigl[\langle\nabla f(w_t),\tilde g_t\rangle\mid w_t\bigr]
   =\|\nabla f(w_t)\|^2 .                                    \tag{6}
\]
[Step 2]

\subsection*{Step 3 – Bound the noisy gradient norm}
From Lemma~\ref{lem:exp} with $a=\tilde g_t$, $b=\xi_t$,
\[
\|g_t\|^2 = \|\tilde g_t\|^2 + 2\langle\tilde g_t,\xi_t\rangle +\|\xi_t\|^2 .
\] 
Taking full expectation and using $\E[\xi_t]=0$,
\[
\E\!\bigl[\|g_t\|^2\bigr]
   = \E\!\bigl[\|\tilde g_t\|^2\bigr] + \E\!\bigl[\|\xi_t\|^2\bigr] . \tag{7}
\] 
Now,
\[
\|\tilde g_t\|^2 = \|\tilde g_t-\nabla f(w_t)+\nabla f(w_t)\|^2
                = \|\nabla f(w_t)\|^2
                  +\|\tilde g_t-\nabla f(w_t)\|^2
                  +2\langle\nabla f(w_t),\tilde g_t-\nabla f(w_t)\rangle .
\] 
The cross term has zero expectation (unbiasedness), thus
\[
\E\!\bigl[\|\tilde g_t\|^2\bigr]
   = \E\!\bigl[\|\nabla f(w_t)\|^2\bigr]
     +\E\!\bigl[\|\tilde g_t-\nabla f(w_t)\|^2\bigr]
   \le \E\!\bigl[\|\nabla f(w_t)\|^2\bigr] + \zeta^{2}.          \tag{8}
\] 
Finally, by Assumption~\ref{ass:noise},
\[
\E\!\bigl[\|\xi_t\|^2\bigr]=d\sigma^{2}.                        \tag{9}
\] 
Combining (8) and (9) in (7) yields
\[
\E\!\bigl[\|g_t\|^2\bigr]
   \le \E\!\bigl[\|\nabla f(w_t)\|^2\bigr] + \zeta^{2}+d\sigma^{2}. \tag{10}
\] 
[Step 3]

\subsection*{Step 4 – Take expectation of the smoothness bound}
Take total expectation of \eqref{4} and substitute the results of
Steps~2 and~3:
\begin{align}
\E[f(w_{t+1})]
&\le \E[f(w_t)]
   -\eta\,\E\!\bigl[\|\nabla f(w_t)\|^2\bigr]
   +\frac{L\eta^{2}}{2}\,
      \E\!\bigl[\|g_t\|^2\bigr]                                     \nonumber\\
&\le \E[f(w_t)]
   -\eta\,\E\!\bigl[\|\nabla f(w_t)\|^2\bigr]
   +\frac{L\eta^{2}}{2}
      \Bigl(\E\!\bigl[\|\nabla f(w_t)\|^2\bigr] +\zeta^{2}+d\sigma^{2}\Bigr). \tag{11}
\end{align}
[Step 4]

\subsection*{Step 5 – Rearrange to isolate the gradient term}
Collect the $\E[\|\nabla f(w_t)\|^2]$ terms:
\[
\E[f(w_{t+1})]
\le \E[f(w_t)]
   -\Bigl(\eta-\frac{L\eta^{2}}{2}\Bigr)
      \E\!\bigl[\|\nabla f(w_t)\|^2\bigr]
   +\frac{L\eta^{2}}{2}\bigl(\zeta^{2}+d\sigma^{2}\bigr).          \tag{12}
\] 
Because the stepsize satisfies $\eta\le 1/L$, we have
$\eta-\frac{L\eta^{2}}{2}\ge \frac{\eta}{2}$.
Thus
\[
\E[f(w_{t+1})]
\le \E[f(w_t)]
   -\frac{\eta}{2}\E\!\bigl[\|\nabla f(w_t)\|^2\bigr]
   +\frac{L\eta^{2}}{2}\bigl(\zeta^{2}+d\sigma^{2}\bigr).          \tag{13}
\] 
[Step 5]

\subsection*{Step 6 – Relate gradient norm to suboptimality}
For convex $f$, the following inequality holds (a standard consequence of
convexity):
\[
f(w_t)-f(w^\star) \le \langle\nabla f(w_t),w_t-w^\star\rangle
                 \le \| \nabla f(w_t)\|\,\|w_t-w^\star\| .      \tag{14}
\] 
Squaring both sides and using $\|w_t-w^\star\|\le D$ (Assumption~\ref{ass:bounded}),
\[
\|\nabla f(w_t)\|^2 \ge \frac{1}{D^{2}}\bigl(f(w_t)-f(w^\star)\bigr)^2 .
\] 
A weaker but convenient linear bound is obtained via Cauchy–Schwarz:
\[
f(w_t)-f(w^\star) \le \langle\nabla f(w_t),w_t-w^\star\rangle
                 \le \|\nabla f(w_t)\|\,D
\;\Longrightarrow\;
\|\nabla f(w_t)\| \ge \frac{f(w_t)-f(w^\star)}{D}.                \tag{15}
\] 
Hence
\[
\E\!\bigl[\|\nabla f(w_t)\|^2\bigr]
   \ge \frac{1}{D^{2}}\E\!\bigl[(f(w_t)-f(w^\star))^{2}\bigr]
   \ge \frac{1}{D^{2}}\bigl(\E[f(w_t)-f(w^\star)]\bigr)^{2},
\] 
where the last inequality follows from Jensen.
Define $\Delta_t:=\E[f(w_t)-f(w^\star)]\ge0$.
Using the linear bound (15) we obtain a simpler inequality:
\[
\E\!\bigl[\|\nabla f(w_t)\|^2\bigr]\ge\frac{\Delta_t^{2}}{D^{2}}. \tag{16}
\] 
[Step 6]

\subsection*{Step 7 – Recursive inequality for the suboptimality}
Insert (16) into (13):
\[
\Delta_{t+1}
\le \Delta_t -\frac{\eta}{2}\frac{\Delta_t^{2}}{D^{2}}
   +\frac{L\eta^{2}}{2}\bigl(\zeta^{2}+d\sigma^{2}\bigr).           \tag{17}
\] 
Rearrange:
\[
\frac{\eta}{2D^{2}}\Delta_t^{2}
\le \Delta_t-\Delta_{t+1}
   +\frac{L\eta^{2}}{2}\bigl(\zeta^{2}+d\sigma^{2}\bigr).           \tag{18}
\] 
Summing \eqref{18} over $t=0,\dots,T-1$ telescopes the first difference:
\[
\frac{\eta}{2D^{2}}\sum_{t=0}^{T-1}\Delta_t^{2}
\le \Delta_0-\Delta_T
   +\frac{TL\eta^{2}}{2}\bigl(\zeta^{2}+d\sigma^{2}\bigr).         \tag{19}
\] 
Since $\Delta_T\ge0$, drop it and use $\Delta_0\le \frac{L}{2}\|w_0-w^\star\|^{2}$
(which follows from smoothness and convexity).  Moreover,
$\Delta_t^{2}\ge \Delta_T^{2}$ for all $t$, therefore
\[
\frac{\eta T}{2D^{2}}\Delta_T^{2}
\le \frac{L}{2}\|w_0-w^\star\|^{2}
   +\frac{TL\eta^{2}}{2}\bigl(\zeta^{2}+d\sigma^{2}\bigr).          \tag{20}
\] 
Divide by $\frac{\eta T}{2D^{2}}$ and take square roots (using
$\sqrt{a+b}\le\sqrt{a}+\sqrt{b}$) to obtain
\[
\Delta_T
\le \frac{D^{2}}{\eta T}\|w_0-w^\star\|
   +\eta L\bigl(\zeta^{2}+d\sigma^{2}\bigr)^{1/2}.
\] 
Because $D\ge\|w_0-w^\star\|$, we can replace $D^{2}$ by $\|w_0-w^\star\|^{2}$
to get a cleaner bound:
\[
\boxed{\;
\E\!\bigl[f(w_T)-f(w^\star)\bigr]
\le
\frac{\|w_0-w^\star\|^{2}}{2\eta T}
+\frac{\eta L}{2}\bigl(\zeta^{2}+d\sigma^{2}\bigr)
\;}.                                                               \tag{21}
\] 
This is exactly the statement of Theorem~\ref{thm:main}.   ∎
[Step 7]

\subsection*{Step 8 – Proof for the averaged iterate}
The same recursion (13) holds for each $t$.  Summing (13) from $0$ to $T-1$
and dividing by $T$ yields
\[
\frac{1}{T}\sum_{t=0}^{T-1}\E\!\bigl[\|\nabla f(w_t)\|^2\bigr]
\le \frac{2}{\eta T}\bigl(f(w_0)-f(w^\star)\bigr)
   +L\eta\bigl(\zeta^{2}+d\sigma^{2}\bigr).                       \tag{22}
\] 
By convexity,
$f(\bar w_T)-f(w^\star)\le \frac{1}{T}\sum_{t=0}^{T-1}\bigl(f(w_t)-f(w^\star)\bigr)$,
and using the inequality $\|\nabla f(w_t)\|^2\ge
\frac{2}{\eta}\bigl(f(w_t)-f(w^\star)\bigr)-L\eta\bigl(\zeta^{2}+d\sigma^{2}\bigr)$
derived from (13) we directly obtain (3).  ∎
[Step 8]

%%%%%%%%%%%%%%%%%%%%%%%%%%%%%%%%%%%%%%%%%%%%%%%%%%%%%%%%%%%%%%%%%%%%%%%%%%%%%%%%
\section*{Rate Derivation (With Constants)}
\addcontentsline{toc}{section}{Rate Derivation (With Constants)}
From \eqref{21} choose a constant stepsize
\[
\eta = \sqrt{\frac{\|w_0-w^\star\|^{2}}
                 {L\bigl(\zeta^{2}+d\sigma^{2}\bigr)T}} .
\] 
Plugging this $\eta$ into \eqref{21} gives
\[
\E\!\bigl[f(w_T)-f(w^\star)\bigr]
\;\le\;
\sqrt{\frac{L\bigl(\zeta^{2}+d\sigma^{2}\bigr)}
           {T}}\;\|w_0-w^\star\|
\;=\;
O\!\Bigl(\frac{1}{\sqrt{T}}\Bigr)
   +O\!\Bigl(\sigma\sqrt{\frac{dL}{T}}\Bigr).
\] 
Thus the utility loss scales as $1/\sqrt{T}$ plus a term proportional to
the noise magnitude required for differential privacy.

%%%%%%%%%%%%%%%%%%%%%%%%%%%%%%%%%%%%%%%%%%%%%%%%%%%%%%%%%%%%%%%%%%%%%%%%%%%%%%%%
\section*{Special Cases}
\addcontentsline{toc}{section}{Special Cases}
\begin{itemize}
\item \textbf{No privacy noise ($\sigma=0$).}  Bound \eqref{21}
reduces to the classical SGD rate
$\frac{\|w_0-w^\star\|^{2}}{2\eta T} + \frac{\eta L}{2}\zeta^{2}$.
\item \textbf{Exact gradients ($\zeta=0$) and full‑batch ($b=n$).}
  The stochastic variance disappears and only the Gaussian term $d\sigma^{2}$
  remains.
\item \textbf{Strongly convex $f$ (parameter $\mu>0$).}
  With the same proof, one obtains a linear convergence term
  $\bigl(1-\eta\mu\bigr)^{T}$ plus the additive noise floor
  $O(\eta\sigma^{2})$.
\end{itemize}

\end{document}